% Historical pre-NoBatchDist text. Excluded from the active manuscript.
% Values retain their old IGD metric and original version identity.
\subsection{Component Ablation of PAQC and CDIS}
\label{sec:exp:ablation}

The ablation compares four configurations at a representative count fixed at
$k=1.5M$, which is 15 at $M=10$ and 30 at $M=20$: REMO under the same
representative count, which provides the reference grouping and selection
rules; the complete framework with CDIS removed (Full-CDIS); the complete
framework with PAQC removed (Full-PAQC); and the complete framework (Full,
i.e.\ PACDIS). Expressing the ablation as removals from the complete framework
keeps the shared grouping, selection and representative rules identical, so
that only the component under test changes. The significance symbols in the
table are the exported comparisons of each variant against REMO.

Table~\ref{tab:exp:ablation} reports the results on all DTLZ and WFG
problems at both objective counts. The decision dimension is 30, except for
WFG2 and WFG3, which use 31.

\begin{table*}[tp]
\centering
\caption{Ablation study of PAQC and CDIS on all DTLZ and WFG test problems with 10 and 20 objectives.}
\label{tab:exp:ablation}
\footnotesize
\setlength{\tabcolsep}{3.5pt}
\renewcommand{\arraystretch}{1.08}
\begin{tabular*}{\textwidth}{@{\extracolsep{\fill}}lccccc@{}}
\toprule
Problem & $D$ & Full-CDIS & Full-PAQC & Full & REMO \\
\midrule
\multicolumn{6}{c}{\textit{10-objective problems}} \\
\midrule
DTLZ1 & 30 & $2.2265e+2\;(2.72e+1)^{+}$ & $2.4443e+2\;(4.93e+1)^{=}$ & $\mathbf{2.1675e+2}\;(3.40e+1)^{+}$ & $2.5572e+2\;(3.72e+1)$ \\
DTLZ2 & 30 & $1.1325e+0\;(1.03e-1)^{=}$ & $\mathbf{1.1205e+0}\;(9.62e-2)^{=}$ & $1.1384e+0\;(8.34e-2)^{=}$ & $1.1321e+0\;(8.14e-2)$ \\
DTLZ3 & 30 & $\mathbf{9.4062e+2}\;(1.36e+2)^{=}$ & $1.0107e+3\;(1.74e+2)^{=}$ & $9.5539e+2\;(8.73e+1)^{=}$ & $9.8158e+2\;(1.46e+2)$ \\
DTLZ4 & 30 & $\mathbf{1.1720e+0}\;(8.09e-2)^{+}$ & $1.2488e+0\;(9.91e-2)^{=}$ & $1.2066e+0\;(1.25e-1)^{=}$ & $1.2536e+0\;(1.07e-1)$ \\
DTLZ5 & 30 & $\mathbf{6.0132e-1}\;(1.35e-1)^{=}$ & $7.0299e-1\;(9.35e-2)^{=}$ & $6.9825e-1\;(8.07e-2)^{=}$ & $6.5860e-1\;(1.07e-1)$ \\
DTLZ6 & 30 & $1.5000e+1\;(5.75e-1)^{+}$ & $1.6175e+1\;(8.26e-1)^{=}$ & $\mathbf{1.4606e+1}\;(7.52e-1)^{+}$ & $1.6101e+1\;(6.03e-1)$ \\
DTLZ7 & 30 & $1.7305e+1\;(1.70e+0)^{+}$ & $1.4415e+1\;(1.84e+0)^{+}$ & $\mathbf{1.0120e+1}\;(1.55e+0)^{+}$ & $2.1381e+1\;(2.48e+0)$ \\
WFG1 & 30 & $3.2000e+0\;(3.80e-2)^{+}$ & $\mathbf{3.1660e+0}\;(5.55e-2)^{+}$ & $3.1698e+0\;(3.68e-2)^{+}$ & $3.2346e+0\;(5.80e-2)$ \\
WFG2 & 31 & $\mathbf{2.1621e+0}\;(6.46e-1)^{+}$ & $2.7797e+0\;(9.10e-1)^{=}$ & $2.4058e+0\;(9.43e-1)^{=}$ & $2.7019e+0\;(8.69e-1)$ \\
WFG3 & 31 & $1.6399e+0\;(5.18e-2)^{-}$ & $1.6242e+0\;(9.57e-2)^{-}$ & $1.7180e+0\;(7.95e-2)^{-}$ & $\mathbf{1.5502e+0}\;(8.35e-2)$ \\
WFG4 & 30 & $7.1851e+0\;(9.13e-1)^{=}$ & $6.8829e+0\;(5.64e-1)^{=}$ & $\mathbf{6.8062e+0}\;(6.20e-1)^{=}$ & $7.2123e+0\;(7.13e-1)$ \\
WFG5 & 30 & $\mathbf{5.8962e+0}\;(3.37e-1)^{=}$ & $5.9358e+0\;(4.73e-1)^{=}$ & $6.0291e+0\;(5.22e-1)^{=}$ & $5.9358e+0\;(4.19e-1)$ \\
WFG6 & 30 & $5.3355e+0\;(3.10e-1)^{=}$ & $5.3791e+0\;(2.61e-1)^{=}$ & $\mathbf{5.2253e+0}\;(2.03e-1)^{+}$ & $5.4737e+0\;(2.95e-1)$ \\
WFG7 & 30 & $5.4193e+0\;(2.07e-1)^{+}$ & $5.6608e+0\;(3.68e-1)^{=}$ & $\mathbf{5.2610e+0}\;(2.42e-1)^{+}$ & $5.7119e+0\;(4.43e-1)$ \\
WFG8 & 30 & $\mathbf{5.6556e+0}\;(1.43e-1)^{+}$ & $5.7850e+0\;(2.65e-1)^{+}$ & $5.7525e+0\;(2.81e-1)^{+}$ & $6.0420e+0\;(3.51e-1)$ \\
WFG9 & 30 & $6.7894e+0\;(6.20e-1)^{=}$ & $6.5815e+0\;(4.11e-1)^{+}$ & $\mathbf{6.4887e+0}\;(5.71e-1)^{+}$ & $6.9496e+0\;(4.45e-1)$ \\
\midrule
\multicolumn{2}{l}{$+/-/=$ vs. REMO} & $8/1/7$ & $4/1/11$ & $8/1/7$ & -- \\
\midrule
\multicolumn{6}{c}{\textit{20-objective problems}} \\
\midrule
DTLZ1 & 30 & $8.1313e+1\;(2.10e+1)^{+}$ & $9.1993e+1\;(2.13e+1)^{=}$ & $\mathbf{7.7516e+1}\;(9.40e+0)^{+}$ & $1.0194e+2\;(2.14e+1)$ \\
DTLZ2 & 30 & $1.1979e+0\;(3.90e-2)^{-}$ & $\mathbf{1.1445e+0}\;(5.21e-2)^{=}$ & $1.1452e+0\;(6.13e-2)^{=}$ & $1.1484e+0\;(4.27e-2)$ \\
DTLZ3 & 30 & $3.2741e+2\;(5.90e+1)^{+}$ & $3.8293e+2\;(7.44e+1)^{=}$ & $\mathbf{3.1754e+2}\;(6.18e+1)^{+}$ & $3.7834e+2\;(8.48e+1)$ \\
DTLZ4 & 30 & $\mathbf{1.0199e+0}\;(4.38e-2)^{+}$ & $1.1130e+0\;(5.08e-2)^{-}$ & $1.0626e+0\;(4.63e-2)^{=}$ & $1.0608e+0\;(3.77e-2)$ \\
DTLZ5 & 30 & $\mathbf{2.6667e-1}\;(4.58e-2)^{+}$ & $3.4665e-1\;(4.32e-2)^{=}$ & $3.0755e-1\;(4.50e-2)^{=}$ & $3.3109e-1\;(4.78e-2)$ \\
DTLZ6 & 30 & $\mathbf{6.5356e+0}\;(7.58e-1)^{+}$ & $7.4297e+0\;(5.46e-1)^{=}$ & $6.6024e+0\;(5.49e-1)^{+}$ & $7.5855e+0\;(4.73e-1)$ \\
DTLZ7 & 30 & $2.5637e+1\;(5.00e+0)^{+}$ & $2.5967e+1\;(4.34e+0)^{+}$ & $\mathbf{1.8855e+1}\;(3.18e+0)^{+}$ & $4.2108e+1\;(7.08e+0)$ \\
WFG1 & 30 & $5.6149e+0\;(7.80e-2)^{=}$ & $5.5954e+0\;(6.36e-2)^{+}$ & $\mathbf{5.5822e+0}\;(7.14e-2)^{+}$ & $5.6481e+0\;(4.89e-2)$ \\
WFG2 & 31 & $5.1016e+0\;(6.88e-1)^{+}$ & $6.0665e+0\;(1.40e+0)^{+}$ & $\mathbf{5.0670e+0}\;(5.79e-1)^{+}$ & $7.0318e+0\;(1.58e+0)$ \\
WFG3 & 31 & $2.6644e+0\;(1.17e-1)^{-}$ & $2.6447e+0\;(1.60e-1)^{-}$ & $2.7603e+0\;(9.00e-2)^{-}$ & $\mathbf{2.3254e+0}\;(1.85e-1)$ \\
WFG4 & 30 & $2.5291e+1\;(1.53e+0)^{=}$ & $2.4509e+1\;(1.47e+0)^{+}$ & $\mathbf{2.4246e+1}\;(1.32e+0)^{+}$ & $2.5899e+1\;(1.45e+0)$ \\
WFG5 & 30 & $2.3849e+1\;(1.38e+0)^{=}$ & $2.3214e+1\;(1.19e+0)^{=}$ & $\mathbf{2.2696e+1}\;(1.09e+0)^{=}$ & $2.3181e+1\;(1.06e+0)$ \\
WFG6 & 30 & $\mathbf{1.9109e+1}\;(1.02e+0)^{+}$ & $2.0258e+1\;(1.05e+0)^{+}$ & $1.9385e+1\;(1.14e+0)^{+}$ & $2.1493e+1\;(7.92e-1)$ \\
WFG7 & 30 & $\mathbf{2.0450e+1}\;(8.73e-1)^{+}$ & $2.1328e+1\;(1.17e+0)^{+}$ & $2.0818e+1\;(1.05e+0)^{+}$ & $2.2442e+1\;(1.43e+0)$ \\
WFG8 & 30 & $2.1098e+1\;(5.56e-1)^{+}$ & $2.1352e+1\;(1.01e+0)^{+}$ & $\mathbf{2.0711e+1}\;(8.06e-1)^{+}$ & $2.2718e+1\;(1.02e+0)$ \\
WFG9 & 30 & $2.4763e+1\;(1.35e+0)^{=}$ & $2.3910e+1\;(1.49e+0)^{=}$ & $\mathbf{2.3568e+1}\;(1.58e+0)^{=}$ & $2.4590e+1\;(1.41e+0)$ \\
\midrule
\multicolumn{2}{l}{$+/-/=$ vs. REMO} & $10/2/4$ & $7/2/7$ & $10/1/5$ & -- \\
\bottomrule
\end{tabular*}
\par\smallskip
\begin{minipage}{\textwidth}\footnotesize
All variants use the same representative count $k=\lceil1.5M\rceil$. \textit{Full-CDIS} denotes the complete framework without CDIS, whereas \textit{Full-PAQC} denotes the complete framework without PAQC. REMO with the same representative count is used as the reference baseline. Each entry reports the mean IGD with the standard deviation in parentheses, and smaller values are better. The best mean among the four variants is highlighted in bold. Symbols $+$, $-$, and $=$ indicate that the corresponding variant is significantly better than, significantly worse than, or statistically indistinguishable from REMO, respectively.
\end{minipage}
\end{table*}


Full has the lowest mean of the four variants on 7 of the 16 problems at
$M=10$ and 9 of the 16 at $M=20$; Full-CDIS on 6 and 5, Full-PAQC on 2 and
1, and REMO on 1 and 1. Measured against REMO, the exported outcomes for
Full are 8/1/7 at $M=10$ and 10/1/5 at $M=20$ (better/worse/no detected
difference), for Full-CDIS 8/1/7 and 10/2/4, and for Full-PAQC 4/1/11 and
7/2/7. The complete framework is therefore better than the reference on
eight of the sixteen problems at $M=10$ and ten at $M=20$, with a single
loss at each objective count.

DTLZ7 shows the largest improvement associated with the complete framework:
relative to REMO, its mean IGD decreases from
$2.1381e+1$ to $1.0120e+1$ at $M=10$ and from $4.2108e+1$ to $1.8855e+1$ at
$M=20$, reductions of 52.7\% and 55.2\%. Most of the change at $M=10$ is
associated with CDIS: relative to Full-CDIS, Full lowers the mean by 41.5\%,
and relative to Full-PAQC by 29.8\%. At $M=20$ the two single-module means
are nearly equal ($2.5637e+1$ and $2.5967e+1$), and Full lowers them by
26.5\% and 27.4\%.

The pattern is not uniform. REMO has the lowest mean on WFG3 at both
objective counts, where all three variants are marked $-$. Removal of PAQC
also degrades DTLZ5 at both objective counts: Full-PAQC is the worst variant
on that problem, with $7.0299e-1$ at $M=10$ and $3.4665e-1$ at $M=20$,
against $6.0132e-1$ and $2.6667e-1$ for Full-CDIS.

% Queue the wide sensitivity table ahead of its discussion for top placement.
\begin{table*}[tp]
\centering
\caption{Sensitivity to the switching probability $p_{\mathrm{mix}}$ at $M=10$ and $M=20$.}
\label{tab:exp:pmix}
\footnotesize
\setlength{\tabcolsep}{3.6pt}
\renewcommand{\arraystretch}{1.08}
\begin{tabular*}{\textwidth}{@{\extracolsep{\fill}}lcccccc@{}}
\toprule
Problem & $D$ & $p_{\mathrm{mix}}=0$ & $0.25$ & $0.50$ & $0.75$ & $1.00$ \\
\midrule
\multicolumn{7}{c}{\textit{10-objective problems}} \\
\midrule
DTLZ2 & 30 & $1.4673e+0\;(1.04e-1)^{-}$ & $1.1980e+0\;(9.10e-2)^{-}$ & $1.1101e+0\;(1.23e-1)$ & $1.1174e+0\;(1.01e-1)^{=}$ & $\mathbf{1.0575e+0}\;(1.27e-1)^{=}$ \\
DTLZ4 & 30 & $1.4452e+0\;(9.75e-2)^{-}$ & $1.2556e+0\;(8.69e-2)^{=}$ & $1.2554e+0\;(1.06e-1)$ & $\mathbf{1.2135e+0}\;(7.44e-2)^{=}$ & $1.2491e+0\;(1.11e-1)^{=}$ \\
DTLZ7 & 30 & $\mathbf{7.1374e+0}\;(1.16e+0)^{+}$ & $8.6532e+0\;(1.65e+0)^{+}$ & $1.0613e+1\;(1.51e+0)$ & $1.2275e+1\;(2.38e+0)^{-}$ & $1.9427e+1\;(2.99e+0)^{-}$ \\
WFG1 & 30 & $3.2195e+0\;(3.59e-2)^{=}$ & $3.1908e+0\;(3.58e-2)^{=}$ & $3.1888e+0\;(5.33e-2)$ & $3.1902e+0\;(6.01e-2)^{=}$ & $\mathbf{3.1852e+0}\;(5.40e-2)^{=}$ \\
WFG3 & 31 & $1.7658e+0\;(5.63e-2)^{-}$ & $1.7316e+0\;(5.24e-2)^{=}$ & $\mathbf{1.6872e+0}\;(7.88e-2)$ & $1.6955e+0\;(6.56e-2)^{=}$ & $1.6929e+0\;(7.01e-2)^{=}$ \\
WFG8 & 30 & $\mathbf{5.5866e+0}\;(2.28e-1)^{=}$ & $5.7157e+0\;(1.98e-1)^{=}$ & $5.7532e+0\;(2.64e-1)$ & $5.8423e+0\;(2.77e-1)^{=}$ & $5.7964e+0\;(3.59e-1)^{=}$ \\
\midrule
\multicolumn{2}{l}{Average rank} & $3.67$ & $3.33$ & $\mathbf{2.33}$ & $3.17$ & $2.50$ \\
\multicolumn{2}{l}{$+/-/=$ vs. $0.50$} & $1/3/2$ & $1/1/4$ & -- & $0/1/5$ & $0/1/5$ \\
\midrule
\multicolumn{7}{c}{\textit{20-objective problems}} \\
\midrule
DTLZ2 & 30 & $1.2616e+0\;(5.18e-2)^{-}$ & $1.1619e+0\;(6.89e-2)^{=}$ & $1.1415e+0\;(6.26e-2)$ & $1.1292e+0\;(5.42e-2)^{=}$ & $\mathbf{1.1111e+0}\;(3.51e-2)^{=}$ \\
DTLZ4 & 30 & $1.0742e+0\;(4.16e-2)^{=}$ & $1.0558e+0\;(5.35e-2)^{=}$ & $\mathbf{1.0506e+0}\;(5.17e-2)$ & $1.0661e+0\;(4.34e-2)^{=}$ & $1.1148e+0\;(5.15e-2)^{-}$ \\
DTLZ7 & 30 & $2.1062e+1\;(8.87e+0)^{=}$ & $\mathbf{1.6810e+1}\;(4.17e+0)^{=}$ & $1.8298e+1\;(3.42e+0)$ & $2.3106e+1\;(4.24e+0)^{-}$ & $3.2237e+1\;(7.51e+0)^{-}$ \\
WFG1 & 30 & $5.6088e+0\;(4.12e-2)^{=}$ & $5.6173e+0\;(6.19e-2)^{=}$ & $\mathbf{5.5870e+0}\;(5.53e-2)$ & $5.5998e+0\;(5.47e-2)^{=}$ & $5.6154e+0\;(8.21e-2)^{=}$ \\
WFG3 & 31 & $2.8724e+0\;(9.74e-2)^{-}$ & $2.7862e+0\;(1.34e-1)^{=}$ & $2.7794e+0\;(9.92e-2)$ & $2.7655e+0\;(9.01e-2)^{=}$ & $\mathbf{2.7013e+0}\;(1.39e-1)^{=}$ \\
WFG8 & 30 & $2.0465e+1\;(9.21e-1)^{=}$ & $\mathbf{2.0362e+1}\;(8.23e-1)^{=}$ & $2.0926e+1\;(8.66e-1)$ & $2.1141e+1\;(1.17e+0)^{=}$ & $2.1744e+1\;(1.34e+0)^{=}$ \\
\midrule
\multicolumn{2}{l}{Average rank} & $3.67$ & $2.83$ & $\mathbf{2.17}$ & $2.83$ & $3.50$ \\
\multicolumn{2}{l}{$+/-/=$ vs. $0.50$} & $0/2/4$ & $0/0/6$ & -- & $0/1/5$ & $0/2/4$ \\
\midrule
\multicolumn{2}{l}{Overall average rank} & $3.67$ & $3.08$ & $\mathbf{2.25}$ & $3.00$ & $3.00$ \\
\multicolumn{2}{l}{Overall $+/-/=$ vs. $0.50$} & $1/5/6$ & $1/1/10$ & -- & $0/2/10$ & $0/3/9$ \\
\bottomrule
\end{tabular*}
\par\smallskip
\begin{minipage}{\textwidth}\footnotesize
Smaller IGD is better, and the lowest mean in each row is highlighted in bold. Symbols $+$, $-$, and $=$ indicate that the corresponding non-default setting is reported as better than, worse than, or not detectably different from $p_{\mathrm{mix}}=0.5$, respectively, following the supplied PlatEMO export. The average-rank rows are computed from the mean IGD values within each objective setting.
\end{minipage}
\end{table*}

Both modules therefore contribute under the reference grouping and selection
rules, and the contribution of each is conditional on the problem. The four
configurations compare every variant with the same reference; they do not
separate an interaction between the two modules from their individual
contributions, and the counts are per-problem comparisons rather than an
omnibus test across problems.

\TODO{Confirm the number of independent runs, the actual terminal FE, the IGD
reference sets and the statistical-test settings from the individual-run
records, and document the selection criterion for the sixteen reported
problems. The exports specify $N=100$ and a configured budget of 300
evaluations; these are configuration values that have not been verified
against the individual runs.}


\subsection{Effect of Selection-Criterion Switching}
\label{sec:exp:pmix}

Table~\ref{tab:exp:pmix} examines whether the default switching probability is a reasonable choice: $p_{\mathrm{mix}}$ was varied over $\{0,0.25,0.5,0.75,1\}$. According to \eqref{eq:mode}, $p_{\mathrm{mix}}=0$ always uses the diversification-oriented criterion when candidate selection is performed, whereas $p_{\mathrm{mix}}=1$ uses the indicator-oriented criterion whenever its surrogate is available; intermediate values stochastically alternate between the two criteria. The exported parameter study contains six benchmark problems at both $M=10$ and $M=20$.



The default $p_{\mathrm{mix}}=0.5$ gives the best overall average rank, $2.25$, compared with $3.67$, $3.08$, $3.00$, and $3.00$ for $p_{\mathrm{mix}}=0$, $0.25$, $0.75$, and $1$, respectively. The same pattern is observed separately at both objective counts: the average rank of the default setting is $2.33$ at $M=10$ and $2.17$ at $M=20$, the lowest among the five settings in each case. This advantage does not arise from uniformly winning every problem. For example, the diversification-only setting performs best on DTLZ7 and WFG8 at $M=10$, while $p_{\mathrm{mix}}=0.25$ obtains the lowest mean on DTLZ7 and WFG8 at $M=20$.

The pairwise annotations nevertheless show a more stable pattern for the equal-probability mixture. Across the 12 problem--objective configurations, $p_{\mathrm{mix}}=0.75$ and $1$ are never reported as better than $p_{\mathrm{mix}}=0.5$, while being worse on two and three configurations, respectively. The $p_{\mathrm{mix}}=0.25$ setting is better once, worse once, and indistinguishable on the remaining ten configurations. The diversification-only setting is better once but worse five times. These results support $p_{\mathrm{mix}}=0.5$ as a robust compromise between the two selection criteria rather than as a universally optimal value on every benchmark instance.
